
\section{Discussion}
\label{sec:discussion}


Several features of the framework are worth highlighting.

\paragraph{Separation of concerns.}
The kernel developer implements the DDM once, exposing $\Mmap$ and its
differential rules as a fixed primitive.  The end user writes the
parameter map~$\Ymap$ in unrestricted host code and obtains
end-to-end gradients automatically.  
Neither party needs to know the internals of the other's contribution.

\paragraph{Generality.}
The decomposition $\Tmap = \Mmap \circ \Ymap$ is not restricted to
static analysis or to any particular element type. 
Any analysis
procedure that can be expressed as a differentiable primitive---linear
or nonlinear statics, transient dynamics, eigenvalue
problems---fits into the same pattern.  
Likewise, the parameter
map~$\Ymap$ can encode arbitrarily complex preprocessing logic:
mesh generation, material homogenisation, stochastic field sampling,
or any combination thereof.

\paragraph{Higher-order composition.}
Because the composed map~$\Tmap$ is itself a differentiable function
in the host framework, it can participate in further compositions.
For instance, an objective function $J(\ModelState)$ defined in host code
yields the fully differentiable pipeline
\begin{equation}\label{eq:full-pipeline}
    \haty \;\xmapsto{\;\Ymap\;}\; \corey
    \;\xmapsto{\;\Mmap\;}\; \ModelState
    \;\xmapsto{\;J\;}\; \mathbb{R},
\end{equation}
whose gradient $\nabla_{\haty} J$ is available via a single reverse-mode pass without any additional implementation.
